\chapter{Annexes}

\section{Arborescence des sources}
\begin{verbatim}
iStudy/
  src/
    main/                 processus principal (Electron + accès PostgreSQL)
      services/           une quarantaine de services métier
      db.js               pool, transactions, LISTEN/NOTIFY
      localization.js     substitution des libellés selon la langue
      labelCatalog.js     résolution des alias non variantes
    preload/index.js      pont contextBridge (API exposée à l'interface)
    renderer/             interface React 19
      components/         écrans et composants partagés
      i18n/               dictionnaires anglais, arabe, français
      index.css           jetons de conception et styles
    shared/               règles communes (mots de passe, semaine, noms)
  deploy/
    server/               scripts SQL et PowerShell du serveur
    client/               préréglage de connexion des postes
    build-offline.ps1     construction des paquets hors ligne
  build/                  icônes et polices embarquées
  scripts/                utilitaires (icônes, migration des médias)
\end{verbatim}

\section{Variables d'environnement du client}
Les variables suivantes remplacent le préréglage de connexion, ce qui permet à un technicien de
corriger un hôte sans modifier de fichier :
\begin{verbatim}
ENOTE_DB_HOST, ENOTE_DB_PORT, ENOTE_DB_NAME,
ENOTE_DB_USER, ENOTE_DB_PASSWORD,
ENOTE_RESOURCE_ROOT, ENOTE_DEPLOYMENT_FILE,
ENOTE_SMB_USER, ENOTE_SMB_PASSWORD
\end{verbatim}

\section{Commandes utiles}
\begin{verbatim}
npm ci                     installation reproductible
npm run dev                lancement en développement
npm run build              construction des trois bundles
npm run build:win          installateur Windows (NSIS + portable)
npm run build:linux        AppImage, deb et archive Linux
npm run migrate:resources  migration des médias vers le partage
\end{verbatim}
Côté serveur :
\begin{verbatim}
psql -U postgres -f deploy/server/01-init-enote.sql
psql -U postgres -d eNote -f deploy/server/02-bootstrap-admin.sql
powershell -File deploy/server/setup-firewall.ps1
powershell -File deploy/server/setup-share.ps1 -SharePath C:\eNoteFiles
powershell -File deploy/server/backup-db.ps1 -BackupDir D:\Backups\eNote
\end{verbatim}

\section{Éléments à compléter}
\begin{itemize}
    \item Identité de l'établissement : sigle, nom complet, ministère de tutelle, direction de
          rattachement et service auteur (page de garde et en-tête) ;
    \item Année de rédaction du document ;
    \item Caractéristiques du poste de développement (modèle, processeur, mémoire, carte graphique) ;
    \item Figures \texttt{images/mpd.pdf} et \texttt{images/usecase.pdf} : modèle physique de données
          et diagramme de cas d'utilisation ;
    \item Captures d'écran des fonctionnalités, nommées comme indiqué sous chaque figure.
\end{itemize}
